% Where the mapping happens, and the two defensible answers.
\begin{tikzpicture}[font=\sffamily\scriptsize, x=1cm, y=1cm]

\node[chlabel, anchor=west] at (-6.6,3.4) {one question: where does a private frame become a standard message};

% ================================================================ option A
\node[layerlabel] at (-6.6,2.9) {A. in the node, which is what this bench does};
\node[fw, text width=24mm] (an) at (-5.4,2.0) {the node\\chapters 1 to 13};
\node[fw, text width=24mm] (am) at (-2.4,2.0) {the mapping\\this chapter};
\node[kern, text width=24mm] (ac) at (0.6,2.0) {the client\\chapter 14};
\node[ext, text width=24mm] (aa) at (3.6,2.0) {the agent};
\node[absentblk, text width=24mm] (ag) at (6.4,2.0) {the graph};
\foreach \a/\b in {an/am, am/ac, ac/aa, aa/ag}{ \draw[flow] (\a) -- (\b); }
\node[chlabel, anchor=north] at (-2.4,1.5) {standard types leave the node};
\node[chlabel, anchor=north] at (3.6,1.5) {and nothing translates them again};

% ================================================================ option B
\node[layerlabel] at (-6.6,0.5) {B. on the host, which is what the maintained stack does};
\node[fw, text width=24mm] (bn) at (-5.4,-0.4) {the node\\its own frames};
\node[hw, text width=24mm] (bb) at (-2.4,-0.4) {the bus\\24-byte frames};
\node[laterblk, text width=24mm] (bm) at (0.6,-0.4) {the master:\\mapping and profile};
\node[ext, text width=24mm] (ba) at (3.6,-0.4) {hardware interfaces};
\node[absentblk, text width=24mm] (bg) at (6.4,-0.4) {the graph};
\foreach \a/\b in {bn/bb, bb/bm, bm/ba, ba/bg}{ \draw[flow] (\a) -- (\b); }
\node[chlabel, anchor=north] at (0.6,-1.25) {this volume does not build this box};
\node[chlabel, anchor=north] at (-2.4,-1.25) {the node stays small};

% ================================================================ the comparison
\node[regcell, fill=usrcol, minimum width=54mm, text width=52mm, inner sep=3pt,
      anchor=north west, align=left] at (-6.6,-2.0)
  {\textbf{A costs} type support, message buffers and about ten times the bytes
   per sample on a constrained part.\\[2pt]
   \textbf{A gives} a node that is a participant in its own right, visible to
   every tool with nothing in between.};
\node[regcell, fill=extcol, minimum width=54mm, text width=52mm, inner sep=3pt,
      anchor=north west, align=left] at (0.4,-2.0)
  {\textbf{B costs} a host component that has to exist, be maintained, and be
   running.\\[2pt]
   \textbf{B gives} a node small enough for a part a manufacturer would put in a
   joint, and a graph that sees a hardware interface.};

\node[umlnote, text width=112mm, anchor=north west] at (-6.6,-4.4)
  {Both are correct and both are shipped. The maintained bus-connected motion
   stack is B, with a driver that implements the standard motion profile and
   exposes it as hardware interfaces, and this volume's node is the device that
   sits opposite it. This bench does A because chapter 14 already put a client on
   the node, and because a volume that never publishes a standard type from a
   microcontroller has not really shown the hard part.};
\end{tikzpicture}
